%======================================================================
\section{The constraint set}
\label{sec:constraints}

\subsection{Normalized units}

Fix one underlying and one expiry $T$.  Everything deterministic ---
discounting, forwards, carry --- divides out before modeling begins.  Let
$F$ be the forward price to $T$ and $D$ the discount factor; we work under
the $T$-forward measure with the \emph{log-forward return} and \emph{gross
return}
\begin{equation}
X=\log\frac{S_T}{F},
\qquad
Y=e^{X}=\frac{S_T}{F} .
\label{eq:Xdef}
\end{equation}
Strikes are quoted in log-moneyness $k=\log(K/F)$, and the
\emph{normalized call} is
\begin{equation}
c(k)=\E\big[(e^{X}-e^{k})^{+}\big],
\qquad
\E[e^{X}]=1 .
\label{eq:calldef}
\end{equation}
The second equality is the martingale (forward) condition: the forward is
the mean of the terminal spot under its own measure.  A dollar quote
$c^{\$}$ converts by $c=c^{\$}/(DF)$, and nothing further about rates or
carry appears in the paper; deterministic carry is a standing assumption,
revisited only when expiries are coupled in \cref{sec:calendar}.  The
normalized put is $P(k)=\E[(e^{k}-e^{X})^{+}]$; put--call parity is
$c(k)-P(k)=1-e^{k}$.

One symbol of time suffices, but it must be the right one.  Desks
distinguish the expiry date, its calendar year fraction, and the
\emph{variance time} --- calendar time possibly dilated around scheduled
events so that implied variance accrues at the market's rate rather than the
clock's.  Everything in this paper that multiplies a variance, scales a
vega, or annualizes a number uses the variance time, written $\tau$; when no
event dilation is in force, $\tau$ is the calendar fraction, and nothing
below depends on which convention feeds it.  Total implied variance at
log-moneyness $k$ is $w(k)=\sigma^{2}(k)\,\tau$, with $\sigma$ the Black
implied volatility defined through the normalized Black call
\begin{equation}
\Black(k,w)=\PhiN(d_{+})-e^{k}\,\PhiN(d_{-}),
\qquad
d_{\pm}=-\frac{k}{\sqrt{w}}\pm\frac{\sqrt{w}}{2},
\label{eq:black}
\end{equation}
by the root $\Black\big(k,w(k)\big)=c(k)$ \citep{Black1976}.  A
\emph{volatility basis point} (vol bp) is $10^{-4}$ in volatility.

\subsection{Statics: what a slice must satisfy}

The reason a smile is not free to be an arbitrary curve is one integration
by parts away.  Write the call as a function of the normalized strike
$y=e^{k}$.  When $Y$ has a density $f_Y$,
\begin{equation}
\frac{\partial c}{\partial y}=-\Prob(Y>y),
\qquad
\frac{\partial^{2}c}{\partial y^{2}}=f_Y(y)\;\ge\;0 .
\label{eq:BL}
\end{equation}
This is the Breeden--Litzenberger identity
\citep{BreedenLitzenberger1978}: the call curve determines the
distribution, and every property of the call curve that a static portfolio
can monetize is a property of the distribution.  A \emph{butterfly} of
half-width $\delta$ --- long one call at $y-\delta$, long one at
$y+\delta$, short two at $y$ --- costs
$c(y-\delta)-2c(y)+c(y+\delta)$, and by \eqref{eq:BL} its price per
$\delta^{2}$ converges to the density at $y$.  A negative butterfly
price is a negative density, exactly.

\begin{definition}[Arbitrage-free slice]
\label{def:slice}
A \emph{slice} is a call curve $c(k)$ of the form \eqref{eq:calldef} for
some law of $X$ with $\E[e^{X}]=1$.  Equivalently, for laws with a density
on $(0,\infty)$: as a function of the strike $y=e^{k}$, $c$ is decreasing
and convex, $(1-y)^{+}\le c(y)\le 1$, the slope lies in $[-1,0]$ with
$c'(y)\to-1$ as $y\to0^{+}$ (anything shallower at zero strike would put
mass at $S_T=0$), and $c(y)\to0$ as $y\to\infty$
\citep{CarrMadan2005}.
\end{definition}

The space of slices is therefore, up to these boundary conditions, the
space of mean-one probability laws on $(0,\infty)$.  That is the object a
fitter must parametrize.  Note one caveat now, because it recurs: convexity
holds in the strike $y$, \emph{not} in log-moneyness $k$.  In $k$ the
chain rule inserts factors of $e^{k}$, and a deep in-the-money call is
legitimately concave in $k$.  Audits that check convexity must check it in
the right variable (\cref{sec:computation}).

\subsection{Why these constraints are awkward in volatility space}

Every trader quotes $\sigma(k)$, so the temptation is to parametrize it
directly; SVI \citep{Gatheral2004,GatheralJacquier2014} is the outstanding
example.  The cost appears when \eqref{eq:BL} is translated into the
volatility chart.  Substituting $c=\Black(k,w(k))$ and differentiating
twice through the Black formula, the density of $X$ factors exactly as
\begin{equation}
f_X(k)\;=\;g_{\mathrm{D}}(k)\,
\frac{\phin\big(d_-(k,w(k))\big)}{\sqrt{w(k)}},
\qquad
g_{\mathrm{D}}(k)=
\Big(1-\frac{k\,w'(k)}{2w(k)}\Big)^{2}
-\frac{w'(k)^{2}}{4}\Big(\frac{1}{w(k)}+\frac14\Big)
+\frac{w''(k)}{2},
\label{eq:durrleman}
\end{equation}
with $d_-$ as in \eqref{eq:black}: no-arbitrage in the volatility chart
is $g_{\mathrm{D}}(k)\ge0$ for all $k$, the condition studied by
Gatheral and Jacquier \citep[\S2]{GatheralJacquier2014}.  The
factorization fixes the normalization of $g_{\mathrm{D}}$, and two
checkpoints let a reader verify it without redoing the two-page
differentiation.  For a flat smile ($w'=w''=0$), $g_{\mathrm{D}}\equiv1$
and \eqref{eq:durrleman} reduces to the Black density of $X$ exactly.
And at the money, substituting the exact derivatives of
\cref{sec:atm} gives $g_{\mathrm{D}}(0)=\sqrt{w_0}\,f_X(0)/\phin(d)\ge0$
--- the identity \eqref{eq:durrleman} read backwards at one point.
Three features make the condition $g_{\mathrm{D}}\ge0$
expensive to live with.  It is \emph{global}: a violation can appear at any
$k$, including between quotes, so it must be monitored on a grid.  It is
\emph{nonlinear} in the parameters, entering through $w$, $w'$, and $w''$
simultaneously, so no useful parametric family satisfies it identically.
And it is \emph{coupled}: the martingale condition and the wing growth
bounds (\cref{sec:tails}) constrain the same parameters from other
directions at the same time.  The volatility chart is the most convenient
to quote and the most expensive to police.  The representation problem of
this paper is to find the chart where the policing cost is zero, and the
next section builds it from first principles.

\subsection{The notation ledger}

The paper insists on few symbols.  \Cref{tab:ledger} lists them all; no
ledger symbol is reused with a second meaning (bound variables local to
a single derivation, such as the butterfly half-width $\delta$ above or
the substitutions inside one proof, stay local and never enter the
ledger).

\begin{table}[htbp]
\centering\small
\setlength{\tabcolsep}{4pt}
\begin{tabular}{@{}llll@{}}
\toprule
$X,\ Y=e^{X}$ & log-forward return; gross return &
$u,\ z=\logit u$ & rank; log-odds \\[2pt]
$k;\ c(k),\ P(k)$ & log-moneyness; norm.\ call, put &
$\Lambda,\ \rho$ & logistic CDF, density \\[2pt]
$Q,\ q=Q'$ & quantile of $X$; quantile density &
$g$ & log transport speed (in rank) \\[2pt]
$x(z),\ m$ & transport map; martingale shift &
$G$ & upper-share ledger \\[2pt]
$w,\ \sigma,\ \tau$ & total variance $\sigma^{2}\tau$; vol; variance time &
$\theta=(L,R,a)$ & coefficients; $P_n$ Legendre \\[2pt]
$\lambda_-,\ \lambda_+$ & endpoint speeds $e^{g(0)},\,e^{g(1)}$ &
$\beta,\ \Psi,\ r$ & Lee slope; Lee map; moment order \\[2pt]
$\Black(k,w)$ & normalized Black call \eqref{eq:black} &
$d=\sqrt{w_0}/2$ & ATM Black argument \\[2pt]
$u_k,\ z_k$ & rank and log-odds of strike $k$ &
$u_*,\ f_*,\ \Delta$ & forward rank; ATM density; digital gap \\[2pt]
$\alpha,\ \nu$ & ridge weight, ridge power &
$\eta,\ h$ & vega floor; haircut half-width \\[2pt]
$\xi$ & logistic-chart coordinate, $\lambda_+=\expit\xi$ &
$N,\ Z_{\max}$ & basis order; quadrature half-width \\
\bottomrule
\end{tabular}
\caption{Every recurring symbol in the paper.  $\PhiN$ and $\phin$ are the
standard normal CDF and PDF; $\expit(z)=1/(1+e^{-z})$ and $\logit$ is its
inverse.  One differentiation convention holds throughout: a \emph{prime}
denotes the derivative of a function of one variable ($c'$, $w'$,
$\sigma'$, $G'$); a \emph{subscripted} $\partial$ denotes a partial
derivative of a function of several variables ($\partial_k\Black$,
$\partial c/\partial\theta_j$).  $Z$ is the logistic score variable of
\cref{sec:model}; $Z_{\max}$, and nothing else, is the quadrature
half-width in log-odds.  $g_{\mathrm{D}}$ \eqref{eq:durrleman} is the
volatility-chart density factor, distinct from the model's log-speed
$g$.}
\label{tab:ledger}
\end{table}
